\subsection{The Section-Average Trace}
\label{sec:trace-cowper}

The section-average trace extracts the $L^2$-optimal rigid approximation 
to the displacement field \eqref{eq:def-section-avg-trace}:
\begin{equation*}
\bar{\bm{u}}(x) := \SectionAverage{\hat{\bm{u}}} - \bar{\bm{\theta}} \times \CentroidLocation, 
\qquad 
\bar{\bm{\theta}}(x) := \RotationAverage{\hat{\bm{u}}},
\end{equation*}
% where $\SectionAverage{\cdot}$ and $\RotationAverage{\cdot}$ are the 
% translation and rotation average functionals of \eqref{eq:def-section-averages}.
where the linear functionals \(\SectionAverage{\cdot}\) and \(\bar{\Theta}[\cdot]\) are defined for \(\bm{v}: [0,L]\times\Section\rightarrow\EuclideanSpace\)
\begin{equation}
\SectionAverage{\bm{u}} := \frac{1}{A}\int_{\Section}\bm{u}\dA,
\qquad
\RotationAverage{\bm u}:=\CentroidStaticInertia^{-1}\!\int_{\Section}(\bm{r}-\CentroidLocation)\times\bm u\dA,
\quad\text{and}\quad
\CentroidStaticInertia:=\int_{\Section}\big(|\CoordinateCentroid|^{2}\mathbf{1}-\CoordinateCentroid\otimes\CoordinateCentroid\big)\dA.
\label{eq:def-section-averages}
\end{equation}
with \(\mathbf{1}\) the 3D identity in \(\mathbb{R}^3\) and \(\CentroidLocation\in\Section\) is the unique point satisfying \(\SectionAverage{\bm{r}-\CentroidLocation} = \mathbf{0}\).

\begin{proposition}[Optimal Rigid Alignment]\label{prop:align-cowper}
The plane-rigid field \(\bar{\bm u}(x)+\bar{\bm\theta}(x)\times\bm{r} \in \PlaneMotions\) that best fits a 3D field \(\hat{\bm u}(x,\bm{r}): [0,L]\times\Section\) in the least‑squares sense is given by:
\begin{equation*}
\bar{\bm u}(\reflenx)=\SectionAverage{\hat{\bm u}}-\bar{\bm\theta}(\reflenx)\times\SectionAverage{\bm{r}},
\qquad
\bar{\bm\theta}(\reflenx)=\bar{\Theta}[\hat{\bm u}],
\end{equation*}
\end{proposition}


\begin{proof}
The section-rigid motion \(\bar{\bm u}+\bar{\bm\theta}\times\bm{r}\) that best fits \(\hat{\bm u}\) in the least‑squares sense minimizes the functional:
\[
\mathcal J(\bar{\bm u},\bar{\bm\theta})
:=
\int_{\Section}\big\| \hat{\bm u}(\bm{r})-(\bar{\bm u}+\bar{\bm\theta}\times\bm{r})\big\|^{2} \dA.
\]

Stationarity of \(\mathcal{J}\) with respect to \(\bar{\bm u}\) gives 
\begin{equation}
\int_{\Section}\!\big(\hat{\bm u}-\bar{\bm u}-\bar{\bm\theta}\times\bm{r}\big)\dA
=\mathbf{0}
\quad\implies\quad
\bar{\bm u}(\reflenx)
=\SectionAverage{\hat{\bm u}} - \bar{\bm\theta}(\reflenx)\times\SectionAverage{\bm{r}}
% =\SectionAverage{\hat{\bm u}}-\bar{\bm\theta}\times\CentroidLocation.
\label{eq:def-avg-ubar}
\end{equation}

Stationarity in \(\bar{\bm\theta}\) requires,
\[
\int_{\Section}\bm{r}\times\big(\hat{\bm u}-\bar{\bm u}-\bar{\bm\theta}\times\bm{r}\big)\dA=\mathbf{0}.
\]
where the identity \(\bm a\cdot(\delta\bar{\bm\theta}\times\bm{r})=(\bm{r}\times\bm a)\cdot\delta\bar{\bm\theta}\) was used.
Insert \(\bar{\bm u}\) from \eqref{eq:def-avg-ubar} and set \(\CoordinateCentroid:=\bm{r}-\CentroidLocation\) so \(\SectionAverage{\CoordinateCentroid}=\mathbf{0}\):
\[
\int_{\Section}\bm{r}\times\big(\hat{\bm u}-\SectionAverage{\hat{\bm u}}\big)\dA
\;-\;
\int_{\Section}\bm{r}\times\big(\bar{\bm\theta}\times(\bm{r}-\CentroidLocation)\big)\dA=\mathbf{0}.
\]
The first term equals \(\int_{\Section}(\bm{r} - \CentroidLocation)\times\hat{\bm u}\dA\).
In the second term, \(\int_{\Section}\bm{r}\times(\bar{\bm\theta}\times\CoordinateCentroid)\dA
=\int_{\Section}\CoordinateCentroid\times(\bar{\bm\theta}\times\CoordinateCentroid)\dA\), and using the vector identity
%$\bm{a} \times(\bm{b} \times \bm{a})=|\bm{a}|^2 \bm{b}-(\bm{a} \cdot \bm{b}) \bm{a}$, 
\(\CentroidLocation \times(\bar{\bm{\theta}} \times \CentroidLocation)=\left(\|\CentroidLocation\|^2 \mathbf{1}-\CentroidLocation \otimes \CentroidLocation\right) \bar{\bm{\theta}}\) we can write:
\[
\int_{\Section}\CoordinateCentroid\times(\bar{\bm\theta}\times\CoordinateCentroid)\dA
=\CentroidStaticInertia\,\bar{\bm\theta},
\]
Hence
\[
\CentroidStaticInertia\,\bar{\bm\theta}
=\int_{\Section}(\bm{r}-\CentroidLocation)\times\hat{\bm u}\dA,\qquad
\bar{\bm\theta}
=\CentroidStaticInertia^{-1}\int_{\Section}(\bm{r}-\CentroidLocation)\times\hat{\bm u} \dA.
\]
where the linear functional \(\bar{\Theta}\) is:
\begin{equation*}
\bar{\Theta}[\bm{u}] := \CentroidStaticInertia^{-1} \int (\bm{r}-\CentroidLocation)\times\bm{u}\dA ,
\qquad 
\CentroidStaticInertia :=\int_{\Omega}\Big(|\bm{r}-\CentroidLocation|^{2}\mathbf{1}-(\bm{r}-\CentroidLocation)\otimes(\bm{r}-\CentroidLocation)\Big)\dA
\end{equation*}
and \(\CentroidStaticInertia\) is the centroidal moment of inertia tensor. 
\end{proof}
Note the properties: 
\begin{subequations}
\begin{align}
\bar{\Theta}[\bm{a}] &= \mathbf{0} 
\qquad \forall \bm{a}\text{ constant in }\bm{r} \label{eq:barTheta-vanish}\\
\bar{\Theta}[\bm{r}] &= \mathbf{0} \\
\bar{\Theta}[\bm{a}\times\bm{r}] &= \bm{a} \label{eq:barTheta-cross-r} \\
\bar{\Theta}\circ\bar{\Theta}[\bm{v}] &= \bar{\Theta}[\bm{v}]
\end{align}
\label{eq:rotavg-identities}
\end{subequations}
\begin{Details}
Furthermore, for any \(\bm v:\Omega\to V\),
\[
\FrameBasis\times\bar\Theta[\FrameBasis\otimes\bm v]
=\IntRGoRG^{-1} \int_{\Omega}(\bm{r}-\CentroidLocation)\otimes\bm v(\bm{r})\dA,
\]
so that
\begin{equation}
\FrameBasis\times\bar\Theta\Big[\FrameBasis\otimes(\bar\Theta[\bm{\WarpShearIesan}]\times(\bm{r}-\CentroidLocation))\Big]
=-\,\bar\Theta[\bm{\WarpShearIesan}]^{\times},
\qquad
\FrameBasis\times\bar\Theta[\FrameBasis\otimes\Section[\bm{\WarpShearIesan}]]
=\Section[\bm{\WarpShearIesan}].
\label{eq:theta-rank-1}
\end{equation}
\[
\RotationAverage{\bm{W}(\bm{r})\bm{a}}
= \CentroidStaticInertia^{-1}\int_\Section \bm{\rho}\times\bm{u}\dA
= \frac{M_\gamma a_\gamma}{J_3}\FrameBasis,
\]
or, more compactly,
\[
\RotationAverage{\bm{W}(\bm{r})\bm{a}}
= \left[\frac{1}{\Io}
\int\big((\bm{r}-\CentroidLocation)\times (\bm{W}(\bm{r})\bm{a})\big)\cdot\mathbf{i} \dA
\right]\FrameBasis
\]
with \(\Io := \operatorname{tr}\CentroidLocation\).
\end{Details}


The trace matrix $\TraceDpDeRoot$ is derived by evaluating the functionals \eqref{eq:def-section-averages} 
to each Saint~Venant mode and extracting the strains at $x = 0$. 
% We then verify that the section-average trace is Class I by confirming that the $x$-dependence 
% follows the universal evolution $\exp\left(x\mathbf{A}\right)$.

\paragraph{Extension Mode.}
The extension displacement is
\begin{equation*}
\hat{\bm{u}}^{(\varepsilon)} = (x\,\FrameBasis - \nu\,\bm{r})\,a_{\varepsilon}.
\end{equation*}
The section averages are:
\begin{align*}
\SectionAverage{\hat{\bm{u}}^{(\varepsilon)}} &= x\,\FrameBasis\,a_\varepsilon - \nu\,\CentroidLocation\,a_\varepsilon, \\
\RotationAverage{\hat{\bm{u}}^{(\varepsilon)}} &= \mathbf{0},
\end{align*}
where use has been made of the rotation average identities \eqref{eq:rotavg-identities}. 
The extracted fields are:
\begin{equation*}
\bar{\bm{u}}^{(\varepsilon)}(x) = x\,\FrameBasis\,a_\varepsilon - \nu\,\CentroidLocation\,a_\varepsilon, 
\qquad 
\bar{\bm{\theta}}^{(\varepsilon)}(x) = \mathbf{0}.
\end{equation*}
The strains are:
\begin{equation*}
\bar{\bm{\gamma}}^{(\varepsilon)} = \FrameBasis\,a_\varepsilon, 
\qquad 
\bar{\bm{\kappa}}^{(\varepsilon)} = \mathbf{0}.
\end{equation*}

\paragraph{Bending Mode.} 
The bending displacement is
\begin{equation*}
\hat{\bm{u}}^{(a)} = \Big( -\tfrac{1}{2}x^2\,\mathbf{1}_\Omega 
+ x\,\FrameBasis \otimes \bm{r} 
+ \bm{W}_{(a)} \Big) \bm{a}_\Omega,
\end{equation*}
where recall $\bm{W}_{(a)} := -\nu\,\mathrm{dev}(\bm{r} \otimes \bm{r})$ is the 
Poisson bending warping \eqref{eq:def-eight-shapes}. 
The section averages give:
\begin{align*}
\bar{\bm{\theta}}^{(a)}(x) &= -x\,(\FrameBasis \times \bm{a}_\Omega) 
+ \RotationAverage{\bm{W}_{(a)}}\,\bm{a}_\Omega.
\end{align*}

The curvature is:
\begin{equation*}
\bar{\bm{\kappa}}^{(a)} = -\FrameBasis \times \bm{a}_\Omega,
\end{equation*}
After simplification using the identity 
$\FrameBasis \times (\FrameBasis \times \bm{a}_\Omega) = -\bm{a}_\Omega$, 
the shear strain is:
\begin{equation*}
\bar{\bm{\gamma}}^{(a)} = \FrameBasis \times \RotationAverage{\bm{W}_{(a)}}\,\bm{a}_\Omega,
\end{equation*}

\paragraph{Torsion Mode.}

The torsion displacement is
\begin{equation*}
\hat{\bm{u}}^{(\vartheta)} = \big( x\,(\FrameBasis \times \bm{r}) 
+ \WarpTwistIesan\,\FrameBasis \big)\,a_\vartheta.
\end{equation*}
The rotation average extracts:
\begin{equation*}
\bar{\bm{\theta}}^{(\vartheta)}(x) = x\,\FrameBasis\,a_\vartheta 
+ \CenterTwistLocation\,a_\vartheta,
\end{equation*}
where the \emph{twist center} $\CenterTwistLocation$ is:
\begin{equation}
\CenterTwistLocation := \RotationAverage{\WarpTwistIesan\,\FrameBasis}.
\label{eq:def-twist-center}
\end{equation}
The center of twist \eqref{eq:def-twist-center} has been attributed to \cite{trefftz1935uber} by \cite{pilkey2003analysis}.

The strains are:
\begin{equation*}
\bar{\varkappa}^{(\vartheta)} = a_\vartheta, 
\qquad 
\bar{\bm{\gamma}}^{(\vartheta)} = (\FrameBasis \times \CenterTwistLocation)\,a_\vartheta,
\end{equation*}

\paragraph{Flexure Mode.}
Recall the particular solution under pure flexure \eqref{eq:iesan-displacement-2}:
\begin{equation*}
\hat{\bm{u}}_{\rm I\!I} = \big( x\,\FrameBasis \times \bm{r} 
+ \WarpTwistIesan\,\FrameBasis \big)\,\CenterShearLocation \cdot \bm{b}_{\Section}
+ \Big( -\tfrac{1}{6}x^3\,\mathbf{1}_\Omega 
+ x\,\PoissonFlexure 
+ \tfrac{1}{2}x^2\,\FrameBasis \otimes (\bm{r} - \CentroidLocation) 
+ \FrameBasis \otimes \bm{\WarpShearIesan} \Big)\,\bm{b}_{\Section},
\end{equation*}
where recall $\PoissonFlexure:\Section\rightarrow\Section\otimes\Section$ is the Poisson flexural warping shape \eqref{eq:def-eight-shapes}.

\subparagraph{Rotation.}
The rotation average \(\RotationAverage{\hat{\bm u}_{\rm I\!I}}\) gives:
\begin{equation*}
\bar{\bm{\theta}}^{(b)}(x) = -\tfrac{1}{2}x^2\,(\FrameBasis \times \bm{b}_{\Section}) 
+ x\,\FrameBasis\otimes\TraceTF\,\bm{b}_{\Section} 
+ \TraceMF\,\bm{b}_{\Section},
% \label{eq:cowper-theta-flexure}
\end{equation*}
where with \(\FrameBasis\otimes\bm{\SectionLocation}_{\pi} := \RotationAverage{\PoissonFlexure}\):
\begin{subequations}
\begin{align}
% \mathbf{K}_1 &:= 
\TraceTF &:= \CenterShearLocation + \RotationAverage{\PoissonFlexure}^{\top}\FrameBasis
=: \CenterShearLocation + \bm{\SectionLocation}_{\pi}, 
\label{eq:cowper-TF} \\
% \TraceTF =: \CenterShearLocation + \FrameBasis^{\top}\RotationAverage{\PoissonFlexure}, 
% \\ %\label{eq:def-K1} \\
%
% \TraceMF 
\bm{\Gamma}_0
&:= \CenterTwistLocation \otimes \CenterShearLocation 
+ \RotationAverage{\FrameBasis \otimes \bm{\WarpShearIesan}}.
\label{eq:cowper-MF}
\end{align}
\end{subequations}
with the center of shear \(\CenterShearLocation\) given by \eqref{eq:def-shear-center}. 
The curvature is:
\begin{equation*}
\bar{\bm{\kappa}}^{(b)}(x) = -x\,(\FrameBasis \times \bm{b}_{\Section}) + \FrameBasis\otimes\TraceTF\,\bm{b}_{\Section}.
% \label{eq:kappa-flexure}
\end{equation*}

% \subparagraph{Shear strain.}
After computing the translation average, displacement derivative, and 
combining with $\FrameBasis \times \bar{\bm{\theta}}^{(b)}$, the quadratic 
terms cancel, yielding:
\begin{equation*}
\bar{\bm{\gamma}}^{(b)}(x) = -x\FrameBasis \otimes \CentroidLocation\,\bm{b}_{\Section} + \TraceFF\,\bm{b}_{\Section},
\end{equation*}
where:
\begin{equation}
\begin{aligned}
\TraceFF 
&:= \SectionAverage{\PoissonFlexure} 
- \CentroidLocation^\times \RotationAverage{\PoissonFlexure}
+ \FrameBasis \times (\CenterTwistLocation \otimes \CenterShearLocation 
+ \RotationAverage{\FrameBasis \otimes \bm{\WarpShearIesan}}),
\end{aligned}
\label{eq:def-Gamma0}
\end{equation}
\begin{Details}
\begin{equation}
\begin{aligned}
\TraceFF &:= (\FrameBasis \times \CentroidLocation) \otimes \CenterShearLocation 
+ \SectionAverage{\PoissonFlexure} 
- \left((\FrameBasis \times \CentroidLocation) \otimes \CenterShearLocation 
+ \CentroidLocation^\times \RotationAverage{\PoissonFlexure}\right)
+ \FrameBasis \times \bm{\Gamma}_0, \\
&:= \SectionAverage{\PoissonFlexure} 
- \CentroidLocation^\times \RotationAverage{\PoissonFlexure}
+ \FrameBasis \times \bm{\Gamma}_0,\\
&:= \SectionAverage{\PoissonFlexure} 
- \CentroidLocation^\times \RotationAverage{\PoissonFlexure}
+ \FrameBasis \times (\CenterTwistLocation \otimes \CenterShearLocation 
+ \RotationAverage{\FrameBasis \otimes \bm{\WarpShearIesan}}),
\end{aligned}
\end{equation}
\end{Details}

\subsubsection{Trace Matrix}

Collecting the contributions at $x = 0$, the trace matrix is:
\begin{equation}
\TraceDeDpRoot = \begin{bmatrix}
1 & \mathbf{0} & 0 & \mathbf{0} %(\TraceFF)_\parallel^\top 
\\[4pt]
\mathbf{0} 
& \mathbf{0} % (\FrameBasis \times \RotationAverage{\bm{W}_{(a)}})_\perp 
& \FrameBasis \times \CenterTwistLocation & \TraceFF \\[4pt]
0 & \mathbf{0} & 1 & \TraceTF^\top %\CenterShearLocation^\top + (\RotationAverage{\PoissonFlexure})_\parallel^\top 
\\[4pt]
\mathbf{0} & -\FrameBasis^\times & \mathbf{0} & \mathbf{0} %\TraceMF
\end{bmatrix}
\qquad\text{and}\qquad 
\TraceDeDpOne = \begin{bmatrix}
0 & \mathbf{0} & 0 & -\CentroidLocation^\top \\[4pt]
\mathbf{0} & \mathbf{0} & \mathbf{0} & \mathbf{0} \\[4pt]
0 & \mathbf{0} & 0 & 0 \\[4pt]
\mathbf{0} & \mathbf{0} & \mathbf{0} & -\FrameBasis^\times
\end{bmatrix}
\label{eq:T0-cowper}
\end{equation}
so \(\TraceBT = \TraceNF = \mathbf{0}\). 
\begin{itemize}
    \item \eqref{eq:req-d1} requires \(\TraceDeDpOne = \TraceDeDpRoot\IesanEvolveSP\), which in view of \eqref{eq:res-d1-block} is satisfied by \eqref{eq:T0-cowper}
    \item \eqref{eq:req-p1} requires \(\TraceDeDpRoot\) exhibit the block form \eqref{eq:T0-plane-structure}, 
\end{itemize}
In view of \eqref{eq:Tinv-from-T}:
\begin{equation}
\begin{array}{rlllll}
% \renewcommand{\arraystretch}{1.3}
\ShiftB & =\left(
\SectionAverage{\PoissonFlexure} 
- \CentroidLocation^\times \RotationAverage{\PoissonFlexure}
+ \FrameBasis \times\RotationAverage{\FrameBasis \otimes \bm{\WarpShearIesan}}
-(\FrameBasis\times\CenterTwistLocation)\otimes \bm{\SectionLocation}_{\pi}
\right)^{-1}
\\
\ShiftK& \!=-\ShiftB\FrameBasis\times\CenterTwistLocation 
\\
\ShiftT & =-\ShiftB(\CenterShearLocation + \bm{\SectionLocation}_{\pi})
\\
\ShiftScaleT & =1 - (\CenterShearLocation + \bm{\SectionLocation}_{\pi})\cdot\ShiftK 
 \\
\ShiftA &=-\FrameBasis^\times (\CenterTwistLocation \otimes \CenterShearLocation + \RotationAverage{\FrameBasis \otimes \bm{\WarpShearIesan}})\,\ShiftB 
& =\mathbf{0}
\\
\ShiftN && = \mathbf{0} 
 \\
\ShiftF & =-\ShiftA(\FrameBasis\times\CenterTwistLocation)
&= \mathbf{0}
 \\
\ShiftScaleN  &=\TraceNF\cdot\ShiftB(\FrameBasis\times\CenterTwistLocation)
& = 0
 \\
\end{array}
\end{equation}

\subsubsection{Summary}

\begin{table}[H]
    \centering
    \caption{Symbols introduced by the geometric trace.}
    % \label{tab:placeholder}
    \begin{tabular}{llcl}
    \toprule
    Symbol & & Units & Description \\ \midrule
    \(\CentroidLocation\)     & \eqref{eq:def-centroid} \\
    \(\CenterShearLocation\)  & \eqref{eq:def-shear-center} \\
    \(\CenterTwistLocation\)  & \eqref{eq:def-twist-center} \\
    \(\bm{\SectionLocation}_{\pi}\) & \(\RotationAverage{\bm{\Pi}_{(b)}}^{\top}\FrameBasis\) \\
    \(\SectionAverage{\bm{\Pi}_{(b)}}\) & \(\tfrac12\Io\oneS - \IntRGoRG\) \\
    \bottomrule
    \end{tabular}
\end{table}

\begin{Details}
\textbf{Verification of Type-3 Exactness}

For the section-average trace to be type-3 exact, the $x$-linear contributions 
must satisfy $\bm{e}_1 = \TraceDeDpRoot\mathbf{A}\,\bm{p}$. 
From the mode-by-mode analysis, the $x$-linear strains arise only from flexure:
\begin{equation*}
\bm{e}_1 = \begin{pmatrix}
(\bm{\Gamma}_1)_\parallel^\top \\
(\bm{\Gamma}_1)_\perp \\
\mathbf{0}^\top \\
-\FrameBasis^\times
\end{pmatrix} \bm{b}_{\Section}
= \begin{pmatrix}
-\CentroidLocation^\top \\
\mathbf{0} \\
\mathbf{0}^\top \\
-\FrameBasis^\times
\end{pmatrix} \bm{b}_{\Section}.
\end{equation*}

Computing $\TraceDeDpRoot\mathbf{A}$: The matrix $\mathbf{A}$ acts only on the 
$\bm{b}_{\Section}$ columns, mapping $\bm{b}_{\Section}$ to 
$(-\CentroidLocation \cdot \bm{b}_{\Section}, \bm{b}_{\Section}, 0, \mathbf{0})^\top$ 
in $(a_\varepsilon, \bm{a}_\Omega, a_\vartheta, \bm{b}_{\Section})$ coordinates. Thus:
\begin{align*}
[\TraceDeDpRoot\mathbf{A}]_{\cdot, \bm{b}_{\Section}} 
&= -\CentroidLocation \cdot [\TraceDeDpRoot]_{\cdot, a_\varepsilon} 
+ [\TraceDeDpRoot]_{\cdot, \bm{a}_\Omega} \\
&= -\CentroidLocation \cdot \begin{pmatrix} 1 \\ \mathbf{0} \\ 0 \\ \mathbf{0} \end{pmatrix}
+ \begin{pmatrix} \mathbf{0}^\top \\ (\FrameBasis \times \RotationAverage{\bm{W}_{(a)}})_\perp \\ \mathbf{0}^\top \\ -\FrameBasis^\times \end{pmatrix} \\
&= \begin{pmatrix} -\CentroidLocation^\top \\ (\FrameBasis \times \RotationAverage{\bm{W}_{(a)}})_\perp \\ \mathbf{0}^\top \\ -\FrameBasis^\times \end{pmatrix}.
\end{align*}

The key observation is that $\RotationAverage{\bm{W}_{(a)}}$ is parallel to 
$\FrameBasis$ (since $\bm{r} \times \bm{W}_{(a)}$ lies in the axial direction), 
so $(\FrameBasis \times \RotationAverage{\bm{W}_{(a)}})_\perp = \mathbf{0}$.

Therefore $\TraceDeDpRoot\mathbf{A} = \bm{e}_1$, confirming 
$\bm{e}(x) = \TraceDeDpRoot \exp\left(x\mathbf{A}\right)\bm{p}$.
\end{Details}

